\documentclass[12pt, a4paper]{article}

% ─── Paquetes ──────────────────────────────────────────────────────────────
\usepackage[utf8]{inputenc}
\usepackage[T1]{fontenc}
\usepackage[spanish]{babel}
\usepackage{geometry}
\usepackage{lmodern}
\usepackage{xcolor}
\usepackage{tcolorbox}
\usepackage{listings}
\usepackage{graphicx}
\usepackage{enumitem}
\usepackage{hyperref}
\usepackage{amsmath}
\usepackage{amssymb}
\usepackage{booktabs}
\usepackage{array}
\usepackage{fancyhdr}
\usepackage{mdframed}
\usepackage{tikz}
\usetikzlibrary{arrows.meta, shapes.geometric, positioning, fit, backgrounds}

\geometry{top=2.5cm, bottom=2.5cm, left=2.8cm, right=2.8cm}

% ─── Colores ───────────────────────────────────────────────────────────────
\definecolor{azulPQC}{RGB}{41, 98, 179}
\definecolor{verdePASS}{RGB}{34, 139, 34}
\definecolor{rojoALERT}{RGB}{180, 50, 30}
\definecolor{grisClaro}{RGB}{245, 247, 250}
\definecolor{azulClaro}{RGB}{219, 234, 254}
\definecolor{amarillo}{RGB}{255, 251, 219}
\definecolor{naranjaFondo}{RGB}{255, 243, 224}
\definecolor{codigoFondo}{RGB}{32, 36, 54}
\definecolor{codigoTexto}{RGB}{220, 220, 200}

% ─── Cabecera / Pie ────────────────────────────────────────────────────────
\pagestyle{fancy}
\fancyhf{}
\fancyhead[L]{\small\textcolor{azulPQC}{\textbf{Verificación Formal — Tamarin Prover}}}
\fancyhead[R]{\small\textcolor{gray}{UMSA — Proyecto de Grado 2026}}
\fancyfoot[C]{\small\thepage}
\renewcommand{\headrulewidth}{0.4pt}

% ─── Cajas temáticas ───────────────────────────────────────────────────────
\tcbuselibrary{most}

\newtcolorbox{cajaConceto}[2][]{
  colback=azulClaro, colframe=azulPQC,
  fonttitle=\bfseries\small, title=#2, #1,
  boxrule=0.7pt, arc=5pt, left=8pt, right=8pt
}

\newtcolorbox{cajaAnalogía}[2][]{
  colback=naranjaFondo, colframe=orange!70!black,
  fonttitle=\bfseries\small, title=\textbf{Analogía: }#2, #1,
  boxrule=0.7pt, arc=5pt, left=8pt, right=8pt
}

\newtcolorbox{cajaPropiedad}[2][]{
  colback=grisClaro, colframe=azulPQC!60,
  fonttitle=\bfseries\small, title=#2, #1,
  boxrule=0.7pt, arc=5pt, left=8pt, right=8pt
}

\newtcolorbox{cajaResultado}[2][]{
  colback=green!8, colframe=verdePASS,
  fonttitle=\bfseries\small, title=\textcolor{verdePASS}{#2}, #1,
  boxrule=0.9pt, arc=5pt, left=8pt, right=8pt
}

% ─── Estilo de código ──────────────────────────────────────────────────────
\lstdefinelanguage{Tamarin}{
  morekeywords={theory, begin, end, builtins, functions, equations,
                restriction, rule, lemma, exists-trace, all-traces,
                let, in, Fr, In, Out, K, All, Ex, not, and, or},
  morestring=[b]",
  morecomment=[l]{//},
  morecomment=[s]{/*}{*/},
  sensitive=true,
}

\lstdefinestyle{tamarinStyle}{
  language=Tamarin,
  backgroundcolor=\color{codigoFondo},
  basicstyle=\ttfamily\small\color{codigoTexto},
  keywordstyle=\color{cyan!80}\bfseries,
  stringstyle=\color{yellow!80},
  commentstyle=\color{green!60}\itshape,
  numbers=left, numberstyle=\tiny\color{gray!50},
  numbersep=8pt,
  breaklines=true, frame=none,
  tabsize=4, keepspaces=true,
  xleftmargin=12pt
}

\lstset{style=tamarinStyle}

% ─── Hiperlinks ────────────────────────────────────────────────────────────
\hypersetup{
  colorlinks=true,
  linkcolor=azulPQC,
  urlcolor=azulPQC,
  pdftitle={Verificación Formal con Tamarin Prover — Protocolo PQC UMSA},
}

% ═══════════════════════════════════════════════════════════════════════════
\begin{document}
% ═══════════════════════════════════════════════════════════════════════════

% ─── Portada ───────────────────────────────────────────────────────────────
\begin{titlepage}
  \centering
  \vspace*{1.5cm}
  {\Large\bfseries\color{azulPQC}
    Universidad Mayor de San Andrés\\[0.3em]
    Carrera de Informática — FCPN}\\[2cm]

  {\Huge\bfseries Verificación Formal de Seguridad\\[0.5em]
   con Tamarin Prover}\\[1em]

  {\large\color{gray}
    Protocolo de Votación Electrónica Post-Cuántica\\
    \textit{(ML-KEM-768 + ML-DSA-65)}}\\[2.5cm]

  \begin{tcolorbox}[colback=azulClaro, colframe=azulPQC, width=0.82\textwidth,
                    boxrule=1pt, arc=8pt]
    \centering\itshape
    ``¿Cómo sabemos que un protocolo de seguridad realmente es seguro\\
    y no simplemente \textbf{parece} seguro?\\[0.5em]
    La respuesta está en la \textbf{verificación formal matemática}.''
  \end{tcolorbox}

  \vfill
  {\large Tesis de Grado — 2026}\\[0.5em]
  {\small Archivo: \texttt{tamarin/protocolo.spthy}}
\end{titlepage}

% ─── Tabla de contenidos ───────────────────────────────────────────────────
\tableofcontents
\newpage

% ═══════════════════════════════════════════════════════════════════════════
\section{¿Qué es la Verificación Formal y por qué la usamos?}
% ═══════════════════════════════════════════════════════════════════════════

\begin{cajaConceto}{Idea Central}
La \textbf{verificación formal} es una técnica matemática que demuestra,
de manera rigurosa y automática, que un protocolo de seguridad cumple ciertas
propiedades (por ejemplo: \textit{nadie puede ver el contenido del voto},
\textit{nadie puede votar dos veces}) bajo \textbf{todas} las posibles
condiciones y ataques concebibles.
\end{cajaConceto}

\medskip

Cuando un ingeniero diseña un protocolo de comunicación segura, lo hace
siguiendo un conjunto de reglas. Sin embargo, existe un problema fundamental:
los \textbf{errores sutiles} en el diseño pueden no ser evidentes aunque el
código funcione correctamente en la práctica. La historia de la criptografía
está llena de protocolos que parecían seguros pero que fueron vulnerados años
después por ataques que nadie había imaginado.

\medskip

La verificación formal ataca este problema de raíz: en lugar de probar el
protocolo con casos de prueba (que podrían no cubrir todos los escenarios),
se \textbf{demuestra matemáticamente} que la propiedad se cumple para
\textit{cualquier} ataque posible.

\begin{cajaAnalogía}{El inspector de puentes}
Imagine que construye un puente. Un ingeniero puede probarlo con un camión de
10 toneladas, luego con uno de 20, de 50\ldots pero un \textbf{verificador
formal} haría algo diferente: demostraría matemáticamente que el puente soporta
\textit{cualquier carga} dentro de las especificaciones, sin necesidad de
probar cada caso individualmente.

En criptografía, el \emph{Tamarin Prover} es ese verificador: demuestra que
el protocolo resiste \textit{cualquier ataque posible} que pueda concebirse
dentro del modelo de seguridad elegido.
\end{cajaAnalogía}

% ═══════════════════════════════════════════════════════════════════════════
\section{¿Qué es Tamarin Prover?}
% ═══════════════════════════════════════════════════════════════════════════

\textbf{Tamarin Prover} es una herramienta de verificación formal de protocolos
criptográficos desarrollada por investigadores de ETH Zürich y otras
instituciones. Permite describir un protocolo como un conjunto de
\textit{reglas de reescritura} y luego demostrar (o refutar) propiedades de
seguridad de forma automática.

\medskip

\begin{table}[h!]
\centering
\begin{tabular}{>{\bfseries}p{3.5cm} p{9cm}}
\toprule
Característica & Descripción \\
\midrule
Modelo de atacante & \textbf{Dolev-Yao}: el atacante controla toda la red.
                     Puede interceptar, modificar, reenviar o eliminar cualquier
                     mensaje, pero \textit{no puede romper las primitivas
                     criptográficas} (firma, cifrado). \\[4pt]
Tipo de demostración & \textbf{Prueba completa}: verifica \textit{todas} las
                       trazas de ejecución posibles, no solo muestras. \\[4pt]
Lenguaje & \texttt{.spthy}: lenguaje de especificación basado en reglas de
           reescritura de multiconjuntos. \\[4pt]
Motor interno & Usa \textit{Maude} como motor de reescritura simbólica. \\
\bottomrule
\end{tabular}
\caption{Características principales de Tamarin Prover.}
\end{table}

\begin{cajaAnalogía}{El juego del atacante más inteligente}
Piense en Tamarin como si pusiéramos al \textbf{atacante más inteligente del
mundo} a intentar romper nuestro protocolo. Este atacante puede:
\begin{itemize}[itemsep=0pt]
  \item Leer todos los mensajes que circulan por la red.
  \item Bloquear, reenviar o modificar mensajes.
  \item Hacerse pasar por cualquier participante.
  \item Combinar mensajes que haya observado de maneras inesperadas.
\end{itemize}
Si Tamarin confirma que el protocolo es seguro, significa que \textit{incluso
ese atacante omnipotente} no puede violar las propiedades que declaramos.
\end{cajaAnalogía}

% ═══════════════════════════════════════════════════════════════════════════
\section{El Modelo de Criptografía Simbólica}
% ═══════════════════════════════════════════════════════════════════════════

Tamarin no trabaja con los números reales de la criptografía (grandes números
primos, matrices de retículos, etc.). En cambio, usa el modelo
\textbf{simbólico}: trata cada operación criptográfica como una ``caja negra
perfecta''.

\medskip

\begin{cajaConceto}{¿Qué significa \textit{caja negra perfecta}?}
En el modelo simbólico:
\begin{itemize}
  \item \texttt{encaps(pk, k)} produce un ciphertext que \textit{nadie} puede
        descifrar sin la clave privada correspondiente.
  \item \texttt{sign(msg, sk)} produce una firma que \textit{nadie} puede
        falsificar sin la clave privada.
  \item \texttt{h(x)} produce un hash que \textit{nadie} puede invertir.
\end{itemize}
Estas suposiciones corresponden exactamente a las propiedades de seguridad
que ML-KEM-768 y ML-DSA-65 garantizan matemáticamente.
\end{cajaConceto}

\medskip

En nuestro archivo \texttt{protocolo.spthy}, declaramos estas operaciones
mediante la sección \texttt{builtins} y \texttt{functions}:

\begin{lstlisting}[caption={Declaración de primitivas criptográficas en Tamarin.}]
builtins: dest-signing, symmetric-encryption

functions:
    encaps/2,       /* ML-KEM-768: encapsula un secreto k con la clave publica */
    decaps/2,       /* ML-KEM-768: desencapsula y recupera el secreto k */
    ek_from_dk/1,   /* deriva la clave publica de encapsulamiento desde la privada */
    h/1             /* SHA3-256: funcion hash para recibos */

equations:
    /* Propiedad de corrección del KEM: decapsulamr un ct valido devuelve el secreto */
    decaps(dk, encaps(ek_from_dk(dk), r)) = r
\end{lstlisting}

La única \textbf{ecuación} definida modela la propiedad fundamental del
mecanismo de encapsulamiento de claves (KEM): si encapsulamos un secreto $r$
con la clave pública derivada de $dk$, luego desencapsulamos con $dk$,
recuperamos $r$.

% ═══════════════════════════════════════════════════════════════════════════
\section{Las Entidades del Protocolo}
% ═══════════════════════════════════════════════════════════════════════════

El protocolo de votación involucra cuatro entidades que interactúan entre sí:

\medskip

\begin{table}[h!]
\centering
\begin{tabular}{>{\bfseries\color{azulPQC}}p{1.5cm}
                >{\bfseries}p{4.5cm}
                p{7.5cm}}
\toprule
Código & Entidad & Responsabilidad \\
\midrule
AR & Autoridad de Registro &
  Verifica la identidad de los votantes, genera sus claves
  criptográficas y los inscribe en el Padrón Electoral.
  Emite credenciales firmadas con ML-DSA-65. \\[4pt]
SV & Servidor de Votación &
  Recepciona las boletas, verifica las firmas digitales,
  previene el doble voto y custodia las boletas cifradas
  en el Ánfora Digital. \\[4pt]
AC & Autoridad de Conteo &
  Custodia exclusivamente la clave de desencapsulamiento
  (\texttt{dk\_SV}). Al cierre electoral, abre el Ánfora
  Digital y totaliza los votos. \\[4pt]
V  & Votante / Ciudadano &
  Selecciona su candidato, cifra su voto con ML-KEM-768,
  firma la boleta con ML-DSA-65 y la envía al SV. \\
\bottomrule
\end{tabular}
\caption{Roles del protocolo de votación electrónica PQC.}
\end{table}

\begin{cajaAnalogía}{La votación tradicional en papel}
En una elección presencial clásica:
\begin{itemize}[itemsep=0pt]
  \item La \textbf{AR} sería el oficial de mesa que verifica su cédula.
  \item El \textbf{SV} sería la urna sellada que guarda los votos.
  \item La \textbf{AC} sería el notario que abre la urna y cuenta los votos.
  \item El \textbf{Votante} sería usted, marcando su papeleta en secreto.
\end{itemize}
La diferencia con el sistema electrónico es que cada uno de estos pasos está
respaldado por criptografía matemáticamente demostrable.
\end{cajaAnalogía}

% ═══════════════════════════════════════════════════════════════════════════
\section{Las Reglas del Modelo: Las 6 Fases del Protocolo}
% ═══════════════════════════════════════════════════════════════════════════

En Tamarin, cada \textbf{regla} (\texttt{rule}) describe una acción atómica
del protocolo. Cada regla tiene la forma:

\begin{center}
\texttt{[ \textit{Premisas} ] --[ \textit{Acciones} ]-> [ \textit{Conclusiones} ]}
\end{center}

Donde:
\begin{itemize}[itemsep=2pt]
  \item \textbf{Premisas}: hechos que deben existir en el estado actual.
  \item \textbf{Acciones}: etiquetas para los lemas de seguridad.
  \item \textbf{Conclusiones}: hechos nuevos que se producen.
\end{itemize}

\subsection{Fase (a) — Configuración de la Autoridad de Registro}

\begin{lstlisting}[caption={Regla: Configuración de la AR con ML-DSA-65.}]
rule Setup_Autoridad_Registro:
    let  pk_AR = pk(~sk_AR)  in
    [ Fr(~sk_AR) ]
  --[ SetupAR(pk_AR) ]->
    [ !ClavePrivadaAR('AR', ~sk_AR)
    , !ClavePublicaAR('AR', pk_AR)
    , Out(pk_AR)             /* La clave publica se publica abiertamente */
    ]
\end{lstlisting}

\begin{cajaConceto}{¿Qué hace esta regla?}
\textbf{Fr({\textasciitilde}sk\_AR)} genera una clave privada fresca (aleatoria)
para la AR. El símbolo \texttt{!} antes de un hecho significa que es
\textbf{persistente}: una vez creado, puede usarse múltiples veces
(a diferencia de los hechos lineales que se consumen una sola vez).

La clave pública \texttt{pk\_AR} se publica en la red (\texttt{Out}):
cualquiera puede verla, lo que permite verificar las firmas de la AR.
\end{cajaConceto}

\subsection{Fase (a) — Configuración del Servidor de Votación}

\begin{lstlisting}[caption={Regla: SV genera claves ML-KEM-768 y ML-DSA-65.}]
rule Setup_Servidor_Votacion:
    let
        ek_SV = ek_from_dk(~dk_SV)   /* Clave publica de encapsulamiento */
        pk_SV = pk(~sk_SV)           /* Clave publica de firma */
    in
    [ Fr(~dk_SV), Fr(~sk_SV) ]
  --[ SetupSV(ek_SV, pk_SV) ]->
    [ !ClaveDesencapsulamientoAC('AC', ~dk_SV)  /* Solo la AC la conoce */
    , !ClavePublicaKEM_SV('SV', ek_SV)
    , !ClavePublicaFirma_SV('SV', pk_SV)
    , Out(ek_SV)                     /* Publicadas para que los votantes las usen */
    , Out(pk_SV)
    ]
\end{lstlisting}

\begin{cajaConceto}{Separación de responsabilidades}
Note que \texttt{{\textasciitilde}dk\_SV} (la clave privada de desencapsulamiento)
se almacena como hecho persistente en \texttt{'AC'}, pero \textbf{nunca se
publica} (\texttt{Out}). Esto modela que el Servidor de Votación no puede
descifrar los votos que custodia: la clave la tiene exclusivamente la
Autoridad de Conteo.
\end{cajaConceto}

\subsection{Fase (b) — Registro del Ciudadano en el Padrón Electoral}

\begin{lstlisting}[caption={Regla: La AR registra al votante y emite su credencial firmada.}]
rule Registro_Ciudadano_Padron:
    let
        pk_V       = pk(~sk_V)
        credencial = <~idV, pk_V, ~ts_reg>       /* Paquete de identidad */
        firma_AR   = sign(credencial, ~sk_AR)    /* Firmada por la AR (ML-DSA-65) */
    in
    [ Fr(~sk_V), Fr(~ts_reg), Fr(~idV)
    , !ClavePrivadaAR('AR', ~sk_AR)
    ]
  --[ RegistroVotante(~idV, pk_V) ]->
    [ !CredencialVotante(~idV, ~sk_V, credencial, firma_AR)
    , !PadronElectoral(~idV, pk_V)
    ]
\end{lstlisting}

\begin{cajaAnalogía}{El carnet de votante}
Esta regla modela la entrega del ``carnet digital de votante''. La
\textbf{credencial} es un documento que contiene el ID del ciudadano y su
clave pública, firmado por la Autoridad de Registro. Esta firma permite que
el Servidor de Votación verifique después que el votante está habilitado,
sin necesidad de contactar a la AR en el momento del sufragio.
\end{cajaAnalogía}

\subsection{Fase (c) — Emisión del Voto en la Mesa de Votación}

\begin{lstlisting}[caption={Regla: El votante cifra y firma su sufragio.}]
rule Emitir_Voto_Mesa:
    let
        ct_kem       = encaps(ek_SV, ~k_efimera)     /* ML-KEM: encapsula secreto */
        voto_cifrado = senc(~m_voto, ~k_efimera)     /* AES-256-GCM: cifra el voto */
        boleta       = <ct_kem, voto_cifrado, credencial, ~ts_voto>
        firma_V      = sign(boleta, ~sk_V)           /* ML-DSA-65: firma la boleta */
    in
    [ Fr(~k_efimera), Fr(~m_voto), Fr(~ts_voto)
    , !CredencialVotante(~idV, ~sk_V, credencial, firma_AR)
    , !ClavePublicaKEM_SV('SV', ek_SV)
    ]
  --[ VotoEmitido(~idV, ~m_voto)
    , VotoEnviado(~idV, boleta, firma_V)
    , SecretoVoto(~m_voto)             /* Marca al voto como secreto para el lema */
    ]->
    [ Out(<boleta, firma_V, firma_AR>) /* Sale a la red (atacante puede verlo) */
    ]
\end{lstlisting}

\begin{cajaAnalogía}{Meter el sobre en el buzón}
Aquí el votante realiza tres pasos simultáneos:
\begin{enumerate}[itemsep=2pt]
  \item \textbf{Genera un secreto efímero} (\texttt{{\textasciitilde}k\_efimera}):
        como una clave única para este voto específico.
  \item \textbf{Cifra el voto} con ese secreto (nadie más puede leerlo).
  \item \textbf{Firma la boleta} con su clave privada (prueba que es él quien votó).
\end{enumerate}
La acción \texttt{SecretoVoto} marca el contenido del voto para que el
lema de confidencialidad pueda monitorearlo.
\end{cajaAnalogía}

\subsection{Fases (d) y (e) — Recepción y Depósito en el Ánfora Digital}

\begin{lstlisting}[caption={Regla: El SV verifica la boleta y la deposita en el Ánfora.}]
rule Recepcion_Y_Deposito_Anfora:
    let
        credencial  = <idV, pk_V, ts_reg>
        boleta      = <ct_kem, voto_cifrado, credencial, ts_voto>
        recibo_hash = h(<boleta, ~ts_recibo>)
    in
    [ In(<boleta, firma_V, firma_AR>)
    , !ClavePublicaAR('AR', pk_AR)
    , !PadronElectoral(idV, pk_V)
    , Fr(~ts_recibo)
    ]
  --[ Eq(verify(firma_AR, credencial, pk_AR), true)  /* Verifica credencial */
    , Eq(verify(firma_V, boleta, pk_V), true)        /* Verifica firma del votante */
    , VotoUnico(idV)                                 /* Anti-duplicado */
    , VotoAceptado(idV, boleta, firma_V)
    , VotoRegistrado(idV, ct_kem, voto_cifrado)
    ]->
    [ AnforaDigital(idV, ct_kem, voto_cifrado, pk_V, credencial, firma_V, ts_voto)
    , Out(recibo_hash)
    ]
\end{lstlisting}

\begin{cajaConceto}{Doble verificación criptográfica}
El Servidor de Votación realiza \textbf{dos} verificaciones de firma:
\begin{itemize}
  \item \texttt{verify(firma\_AR, credencial, pk\_AR)}: verifica que la
        credencial es auténtica (emitida por la AR legítima).
  \item \texttt{verify(firma\_V, boleta, pk\_V)}: verifica que la boleta
        fue construida y firmada por el votante habilitado.
\end{itemize}
La restricción \texttt{VotoUnico(idV)} garantiza que ningún \texttt{idV}
puede disparar esta regla más de una vez (doble voto imposible).
\end{cajaConceto}

\subsection{Fase (f) — Apertura del Ánfora y Escrutinio Oficial}

\begin{lstlisting}[caption={Regla: La AC desencapsula y descifra cada voto.}]
rule Apertura_Anfora_Y_Escrutinio:
    let
        k_recuperada = decaps(~dk_SV, ct_kem)
        voto_plano   = sdec(voto_cifrado, k_recuperada)
    in
    [ AnforaDigital(idV, ct_kem, voto_cifrado, pk_V, credencial, firma_V, ts_voto)
    , !ClaveDesencapsulamientoAC('AC', ~dk_SV)
    ]
  --[ VotoContado(idV, voto_plano) ]->
    [ Out(<idV, voto_plano>) ]
\end{lstlisting}

\begin{cajaConceto}{¿Por qué solo la AC puede descifrar?}
La clave \texttt{{\textasciitilde}dk\_SV} es persistente con el agente
\texttt{'AC'}: solo la Autoridad de Conteo tiene acceso a ella. Sin esta
clave, el ciphertext \texttt{ct\_kem} es computacionalmente indistinguible
de valores aleatorios (seguridad IND-CCA2 de ML-KEM-768).
\end{cajaConceto}

% ═══════════════════════════════════════════════════════════════════════════
\section{Los Lemas de Seguridad: Las Propiedades a Demostrar}
% ═══════════════════════════════════════════════════════════════════════════

Los \textbf{lemas} son las propiedades de seguridad que Tamarin debe
demostrar matemáticamente. Son el corazón de la verificación formal.

\subsection{Lema 1: Ejecutabilidad (Sanity Check)}

\begin{lstlisting}[caption={Lema de ejecutabilidad del protocolo.}]
lemma ejecutabilidad:
    exists-trace
    "Ex idV m blt sig #t1 #t2 #t3.
        VotoEmitido(idV, m)          @#t1   /* Momento en que se emite */
      & VotoAceptado(idV, blt, sig)  @#t2   /* Momento en que el SV acepta */
      & VotoContado(idV, m)          @#t3   /* Momento en que la AC cuenta */
      & #t1 < #t2                           /* El orden es correcto */
      & #t2 < #t3"
\end{lstlisting}

\begin{cajaConceto}{¿Para qué sirve este lema?}
Este es el ``lema de cordura'' (\textit{sanity check}). Demuestra que
el protocolo \textbf{tiene al menos una traza de ejecución válida}: que
existe una manera de que todo funcione correctamente en orden. Si este lema
falla, significa que el modelo está mal construido y el protocolo nunca
puede completarse. Es análogo a verificar que un avión puede despegar
antes de demostrar que puede aterrizar de forma segura.
\end{cajaConceto}

\subsection{Lema 2: Confidencialidad del Voto (Secreto)}

\begin{lstlisting}[caption={Lema: El atacante nunca aprende el contenido del voto.}]
lemma secreto_voto:
    all-traces
    "All m #t.
        SecretoVoto(m) @#t
        ==>
        not(Ex #t2. K(m) @#t2)"
\end{lstlisting}

\begin{cajaConceto}{Interpretación en lenguaje llano}
Este lema dice: \textbf{``Para todo voto $m$ que fue marcado como secreto
en cualquier traza de ejecución, el conocimiento del atacante nunca incluye
$m$''.}

En Tamarin, \texttt{K(m)} significa ``el atacante conoce $m$''. La negación
\texttt{not(Ex \#t2. K(m) @\#t2)} significa ``no existe ningún momento en
el tiempo en que el atacante aprende $m$''.

Esto modela la propiedad de \textbf{confidencialidad del sufragio}: aunque
el atacante controla la red y puede interceptar todos los mensajes, nunca
puede descubrir por quién votó el ciudadano.
\end{cajaConceto}

\subsection{Lema 3: Integridad del Sufragio}

\begin{lstlisting}[caption={Lema: Todo voto contado fue legítimamente emitido.}]
lemma integridad_voto:
    all-traces
    "All idV m #tc.
        VotoContado(idV, m) @#tc
        ==>
        Ex #te. VotoEmitido(idV, m) @#te & #te < #tc"
\end{lstlisting}

\begin{cajaConceto}{Interpretación en lenguaje llano}
\textbf{``Para todo voto $m$ del ciudadano $idV$ que fue contado en el
escrutinio, debe existir un momento anterior en que ese mismo ciudadano
emitió ese mismo voto $m$.''}

Esta propiedad garantiza que:
\begin{itemize}[itemsep=2pt]
  \item No pueden aparecer votos de la nada (votos fantasma).
  \item El contenido del voto no puede ser alterado en tránsito.
  \item El escrutinio refleja exactamente lo que los ciudadanos decidieron.
\end{itemize}
Corresponde al concepto de \textbf{correspondencia} (\textit{correspondence})
en la literatura de verificación formal de protocolos.
\end{cajaConceto}

\subsection{Lema 4: Autenticidad del Votante (Injective Agreement)}

\begin{lstlisting}[caption={Lema: Cada aceptación corresponde a un envío único.}]
lemma autenticidad_votante:
    all-traces
    "All idV blt sig #ta.
        VotoAceptado(idV, blt, sig) @#ta
        ==>
        ( Ex #ts.
            VotoEnviado(idV, blt, sig) @#ts
          & #ts < #ta
          & not( Ex #ta2.
                VotoAceptado(idV, blt, sig) @#ta2
              & not(#ta = #ta2) ) )"
\end{lstlisting}

\begin{cajaConceto}{Interpretación en lenguaje llano}
Este lema demuestra dos propiedades en una sola fórmula:
\begin{enumerate}
  \item \textbf{Autenticidad}: Cada boleta aceptada por el SV fue
        realmente enviada por el votante legítimo (no es una falsificación).
  \item \textbf{Unicidad (anti-replay)}: La misma boleta no puede ser
        aceptada \textbf{dos veces}. El \texttt{not(Ex \#ta2...)} garantiza
        que no existe un segundo momento de aceptación distinto.
\end{enumerate}
Esto previene los \textbf{ataques de repetición} (\textit{replay attacks}),
donde un atacante intercepta un voto válido y lo reenvía múltiples veces.
\end{cajaConceto}

% ═══════════════════════════════════════════════════════════════════════════
\section{Las Restricciones (Axiomas del Sistema)}
% ═══════════════════════════════════════════════════════════════════════════

Las restricciones son condiciones que el protocolo impone como invariantes
del sistema, independientemente del atacante.

\begin{lstlisting}[caption={Restricciones del modelo.}]
restriction Unicidad_Voto:
    "All id #i #j. VotoUnico(id) @i & VotoUnico(id) @j ==> #i = #j"

restriction Verificacion_Exitosa:
    "All x y #i. Eq(x, y) @i ==> x = y"
\end{lstlisting}

\begin{cajaPropiedad}{Restricción 1: Unicidad del Voto}
Esta restricción dice: \textit{``Si la acción \texttt{VotoUnico(id)} ocurre
en dos momentos distintos $i$ y $j$, entonces esos momentos deben ser el
mismo''.} En otras palabras, la acción \texttt{VotoUnico(id)} solo puede
ocurrir una vez por votante. Esta restricción se usa en la regla de
recepción para garantizar el \textbf{principio de doble voto imposible}.
\end{cajaPropiedad}

\begin{cajaPropiedad}{Restricción 2: Verificación Exitosa}
Esta restricción modela que cuando el servidor \textit{verifica} una firma
(acción \texttt{Eq(verify(...), true)}), la verificación \textbf{realmente}
fue exitosa. Sin esta restricción, Tamarin consideraría trazas donde la
verificación dice ``falso'' pero el protocolo continúa, lo cual no
corresponde a una implementación correcta.
\end{cajaPropiedad}

% ═══════════════════════════════════════════════════════════════════════════
\section{Resultados de la Verificación Formal}
% ═══════════════════════════════════════════════════════════════════════════

Tras ejecutar el comando:
\begin{center}
  \texttt{tamarin-prover tamarin/protocolo.spthy --prove}
\end{center}

Los cuatro lemas de seguridad fueron verificados exitosamente:

\begin{cajaResultado}{Resultados de Verificación — 4/4 Lemas VERIFICADOS}
\begin{table}[h!]
\centering
\begin{tabular}{>{\ttfamily}p{4.5cm} p{5cm} >{\bfseries\color{verdePASS}}p{2cm}}
\toprule
\normalfont Lema & Propiedad Garantizada & Estado \\
\midrule
ejecutabilidad & Traza completa ejecutable & VERIFIED \\[4pt]
secreto\_voto & Confidencialidad del sufragio & VERIFIED \\[4pt]
integridad\_voto & Correspondencia E2E & VERIFIED \\[4pt]
autenticidad\_votante & Anti-replay (Injective Agree.) & VERIFIED \\
\bottomrule
\end{tabular}
\end{table}
\end{cajaResultado}

\medskip

\begin{cajaConceto}{¿Qué significa VERIFIED?}
Cuando Tamarin devuelve \texttt{verified} para un lema, significa que ha
explorado \textbf{todas las trazas de ejecución posibles} del protocolo y
no ha encontrado ni una sola en la que la propiedad sea violada.

Esto es una \textbf{demostración matemática}, no una prueba empírica. El
resultado es válido para cualquier número de votantes, cualquier número de
sesiones y cualquier estrategia que pueda seguir un atacante dentro del
modelo Dolev-Yao.
\end{cajaConceto}

% ═══════════════════════════════════════════════════════════════════════════
\section{Limitaciones del Modelo}
% ═══════════════════════════════════════════════════════════════════════════

Es importante comprender qué garantiza y qué no garantiza esta verificación
formal:

\medskip

\begin{table}[h!]
\centering
\begin{tabular}{>{\bfseries}p{6cm} p{7.5cm}}
\toprule
Lo que SÍ garantiza & Lo que NO cubre este modelo \\
\midrule
Las 4 propiedades de seguridad del protocolo en el modelo simbólico
& La seguridad matemática computacional de los algoritmos (ML-KEM, ML-DSA) \\[4pt]
Resistencia a cualquier ataque de red (Dolev-Yao)
& Ataques de canal lateral (tiempo, consumo de energía) \\[4pt]
Imposibilidad de doble voto en el modelo
& Errores en la implementación en código Python/C \\[4pt]
Autenticidad y confidencialidad simbólica
& Ataques al sistema operativo o hardware del votante \\
\bottomrule
\end{tabular}
\caption{Alcance y limitaciones del modelo formal.}
\end{table}

\medskip

\noindent La verificación formal es un \textbf{complemento} a la implementación
segura, a las pruebas unitarias y a las auditorías de código. No las reemplaza.

% ═══════════════════════════════════════════════════════════════════════════
\section{Cómo Ejecutar la Verificación}
% ═══════════════════════════════════════════════════════════════════════════

Para reproducir la verificación formal en su máquina:

\begin{lstlisting}[language=bash, backgroundcolor=\color{codigoFondo},
                   basicstyle=\ttfamily\small\color{codigoTexto}, numbers=none,
                   caption={Comandos para ejecutar la verificación.}]
# Verificacion automatica de todos los lemas
tamarin-prover tamarin/protocolo.spthy --prove

# Modo interactivo con interfaz web (GUI en navegador)
tamarin-prover interactive tamarin/protocolo.spthy

# Verificar un solo lema especifico
tamarin-prover tamarin/protocolo.spthy --prove=secreto_voto
\end{lstlisting}

\bigskip

\begin{tcolorbox}[colback=green!8, colframe=verdePASS, boxrule=0.9pt, arc=5pt]
\centering\bfseries\color{verdePASS}
Verificación Formal Completada Exitosamente\\[0.3em]
\normalfont\small
4 de 4 lemas verificados $\cdot$ Modelo Dolev-Yao $\cdot$ ML-KEM-768 + ML-DSA-65 (FIPS 203/204)\\
Universidad Mayor de San Andrés — Proyecto de Grado 2026
\end{tcolorbox}

% ─── Bibliografía ──────────────────────────────────────────────────────────
\newpage
\begin{thebibliography}{9}

\bibitem{tamarin}
  Schmidt, B., Meier, S., Cremers, C., \& Basin, D. (2012).
  \textit{The TAMARIN Prover for the Symbolic Analysis of Security Protocols}.
  In CAV 2012, LNCS 7358, pp. 696–701. Springer.

\bibitem{fips203}
  NIST. (2024).
  \textit{FIPS 203 — Module-Lattice-Based Key-Encapsulation Mechanism Standard}.
  National Institute of Standards and Technology.

\bibitem{fips204}
  NIST. (2024).
  \textit{FIPS 204 — Module-Lattice-Based Digital Signature Standard}.
  National Institute of Standards and Technology.

\bibitem{dolevyao}
  Dolev, D., \& Yao, A. (1983).
  \textit{On the security of public key protocols}.
  IEEE Transactions on Information Theory, 29(2), 198–208.

\bibitem{pqclean}
  PQClean Project.
  \textit{Clean, portable, tested implementations of post-quantum cryptography}.
  \url{https://github.com/PQClean/PQClean}

\end{thebibliography}

\end{document}
